\documentclass[12pt,addpoints,answers]{exam}

\usepackage{mystyle}

\title{MATH 3191-001 Fall 2026 Exam 1}
\date{23 September 2026}

\begin{document}
\maketitle

    \paragraph{Instructions:}

    Write your name on the paper and show all work that leads to your answers. Answers
    without work will not be given credit, so you MUST justify your answer. You
    have 1 hour and 15 minutes to take the exam, and there is extra scratch paper available at
    the front of the room. If you need to use it, please label which problem you are working on and
    turn it in with the exam, so I can look at your work if necessary.

    This exam has $115$ points available, but will be scored out of $105$, so a score of $\frac{115}{105}
    \approx 109.5$ is possible.
    \vfill
    \begin{center}
        \gradetable[h]
    \end{center}
    \newpage
\begin{questions}
    \section*{Short Answer Questions}
    \question (20 points total)
    This problem is walking you through a proof of Invertible Matrix Theorem statement 3 implying
    statement 6.
    \begin{parts}
        \part[10] Let $A\in\R^{n\times n}$ be row equivalent to an upper triangular matrix $B$.
        IE there exist elementary matrices $E_1,\dots,E_k$ such that $E_k\cdots E_1\cdot A = B$.
        Argue why if $B$ is invertible, then so is $A$.

        \emph{Hint:} If you're stuck, think about the determinant of the product of matrices and if
        the determinant of each elementary matrix can be $0$.
        \begin{solution}
            We recall that an elementary matrix is a matrix that is $1$ row operation away from the identity.
            This means that for $\ell=1,\dots,k$, the diagonal of $E_\ell$ is non-zero and is
            upper and/or lower triangular. Thus, for $\ell=1,\dots,k$, we know that
            $\det{E_\ell}\neq 0$. From the
            properties of the determinant, we know that
            \begin{eqnarray*}
                \det{E_k\cdots E_1\cdot A} &=& \det{B} \\
                \det{E_k}\cdots\det{E_1}\cdot\det{A} &=& \det{B} \\
                \det{A} &=& \dfrac{\det{B}}{\det{E_k}\cdots\det{E_1}}
            \end{eqnarray*}
            If $B$ is invertible, we know that $\det{B}\neq 0$, so from above we know $\det{A}\neq 0$,
            which lets us conclude that $A$ is invertible as well.
        \end{solution}
        \vfill
        \part[10] Prove that if an upper triangular matrix $A\in\R^{n\times n}$ has a pivot in
        every column, then show the columns of $A$ are linearly independent

        \emph{Hint:} Recall what having $n$ pivots means about the number of solutions to a given system!
        \begin{solution}
            Recall that for a matrix $A$ to have $n$ pivots then the diagonal elements of it
            must be non-zero. This means the augmented matrix $\aMat{c|c}{A & \vec{0}}$ can be row
            reduced to $\aMat{c|c}{I & \vec{0}}$. Thus, the homogeneous system has only the trivial
            solution. In other words if we consider $A$ as $A = \bMat{\vec{a}_1 & \dots & \vec{a}_n}$
            then the system $c_1\vec{a}_1 + \cdots + c_n\vec{a}_n$ has only the solution
            $(c_1,\cdots,c_n) = (0,\cdots,0)$ which, by the definition of linear independence means
            that $\vec{a}_1,\cdots,\vec{a}_n$ are linearly independent.
        \end{solution}
        \vfill
        \note{Combining part (a) and (b) would let us say that any matrix $A\in\R^{n\times n}$ being
        row equivalent to an upper triangular matrix with $n$ pivots has linearly independent columns.}
    \end{parts}
    \newpage
    \question[15] Let $T:\R^{n}\rightarrow\R^{m}$ be a linear transformation, describe what information
    you need to determine if $T$ is invertible

    \emph{Hint:} think about the invertible matrix theorem and a condition on $m$ and $n$.
    \begin{solution}
        5 points each for (1) requiring $m=n$, (2) alluding to the invertible matrix theorem, and
        (3) how $T$ acts on $n$ linearly independent vectors (will give credit if linearly independent is
        not explicitly stated)

        Sample solution:

        First, we need to know that $m=n$, and then we need to know how $T$ acts on $n$ linearly
        independent vectors. Then with this information, we can construct the matrix associated with
        this linear transformation and check if $A$ is row-equivalent to an upper triangular matrix
        with $n$ pivots. If all of this is true, then we get to say that $T$ is an invertible linear
        transformation. Otherwise, we would say that $T$ is not an invertible linear transformation.
    \end{solution}
    \newpage
    \section*{Computational Question}
    \question[15]
    Let $A\in\R^{2\times 3},\vec{b}\in\R^2$ be as given below. Determine the solution set that
    solves the equation $A\vec{x} = \vec{b}$.
    \[
        A = \bMat{
        1 & 4 & 2 \\
        7 & 30 & 12
        } \qquad \vec{b} = \bMat{ 6 \\ 44 }
    \]
    \emph{Hint: your final answer should contain only integers}
    \begin{solution}
        We will construct the following augmented matrix
        \[
            \aMat{c|c}{A & \vec{b}} = \aMat{ccc|c}{
                1 & 4 & 2 & 6 \\
                7 & 30 & 12 & 44
            }
        \]
        Which we row reduce as follows

        \begin{eqnarray*}
            \aMat{ccc|c}{
                1 & 4 & 2 & 6 \\
                7 & 30 & 12 & 44
            } &\xrightarrow{R_2\gets R_2 - 7R_1}&\aMat{ccc|c}{
                1 & 4 & 2 & 6 \\
                0 & 2 & -2 & 2
            } \\
            &\xrightarrow{R_2\gets \frac{1}{2}R_2}&\aMat{ccc|c}{
                1 & 4 & 2 & 6 \\
                0 & 1 & -1 & 1
            }\\
            &\xrightarrow{R_1\gets R_1 - 4R_2}&\aMat{ccc|c}{
                1 & 0 & 6 & 2 \\
                0 & 1 & -1 & 1
            }
        \end{eqnarray*}
        This lets us construct our solution set as
        \[
            \set{\vec{x}\in\R^{3}}{\vec{x} = \bMat{2\\1\\0} + t\bMat{-6\\1\\1}\text{ for some }t\in\R}
        \]
    \end{solution}
    \newpage
    \question[20]
    Let $A\in\R^{3\times 2}, \vec{b}\in\R^{3}$ be as given below, determine a value $z\in\R$
    such that $A\vec{x} = \vec{b}$ is consistent (IE: has at least $1$ solution)
    \[
        A = \bMat{
        2 & z \\
        4 & 4 \\
        6 & 2
        }\qquad \bMat{ 1\\ 4\\ 4}
    \]
    \begin{solution}
        We will first set up the augmented system:
        \[\aMat{c|c}{A & \vec{b}} = \aMat{cc|c}{
            2 & z & 1 \\
            4 & 4 & 4 \\
            6 & 2 & 4
        }\]

        which we row reduce as follows
        \begin{eqnarray*}
            \aMat{cc|c}{
            2 & z & 1 \\
            4 & 4 & 4 \\
            6 & 2 & 4
        } &\xrightarrow[R_2\gets R_2 - 2R_1]{R_3\gets R_3 - 3R_1}&\aMat{cc|c}{
            2 & z & 1 \\
            0 & 4-2z & 2 \\
            0 & 2-3z & 1
        }
        \end{eqnarray*}
        Now, we want this system to be consistent. Since we only have two variables, we must be able
        to get the final row to be a $0$ row. Looking at the last column, the only way this can happen
        is if $R_3\gets R_3 - 2R_2$ results in the entire row being $0$. In other words: we need $z$ to
        satisfy:
        \[
            4 - 2z = 4 - 6z \rightarrow 4z = 0
        \]
        This means that $z$ must be $0$.

        There are other methods that we can do as well!
    \end{solution}
    \newpage
    \question[15]
    Let $T:\R^2\rightarrow\R^3$ be a linear transformation as given below. Determine the matrix
    associated $A$ associated with $T$ ie: compute $A$ such that $A\vec{x} = T(\vec{x})$.

    \[
        T\left(\bMat{1\\3}\right) = \bMat{1\\6\\3} \qquad T\left(\bMat{2\\5}\right) = \bMat{7\\1\\3}
    \]
    \begin{solution}
        We recall that the matrix associated with a linear transformation is defined to be
        \[
            A = \bMat{ T(\vec{e}_1) & \dots & T(\vec{e}_n) }.
        \]

        This means we need to be able to write $\vec{e}_1 = \bMat{1\\0}$ and $\vec{e}_2 = \bMat{0\\1}$
        as a linear combination of $\vec{v}_1 = \bMat{1\\3}$ and $\vec{v}_2 = \bMat{2\\5}$. We will
        do this as we discussed on Monday in class; by row reducing the augmented matrix given by
        \[
            \aMat{cc|c}{ \vec{v}_1 & \vec{v}_2 & I } = \aMat{cc|cc}{
                1 & 2 & 1 & 0 \\
                3 & 5 & 0 & 1
            }.
        \]
        Which we do as follows:
        \begin{eqnarray*}
            \aMat{cc|cc}{
                1 & 2 & 1 & 0 \\
                3 & 5 & 0 & 1
            } &\xrightarrow{R_2\gets R_2 - 3R_1}& \aMat{cc|cc}{
                1 & 2 & 1 & 0 \\
                0 & -1 & -3 & 1
            }\\
            &\xrightarrow{R_2\gets -R_2}&\aMat{cc|cc}{
                1 & 2 & 1 & 0 \\
                0 & 1 & 3 & -1
            } \\
            &\xrightarrow{R_1\gets R_1 - 2R_2}&\aMat{cc|cc}{
                1 & 0 & -5 & 2 \\
                0 & 1 & 3 & -1
            }
        \end{eqnarray*}
        This means that we can write
        $\vec{e}_1 = -5\vec{v}_1 + 3\vec{v}_2$ and $\vec{e}_2 = 2\vec{v}_1 - \vec{v}_2$. Thus we get
        \begin{align*}
            T(\vec{e}_1) &= -5T(\vec{v}_1) + 3T(\vec{v}_2) \\
            &= -5\bMat{1\\6\\3} + 3\bMat{7\\1\\3} \\
            &= \bMat{16\\-27\\-6}\\
            T(\vec{e}_2) &= 2T(\vec{v}_1) - T(\vec{v}_2) \\
            &= 2\bMat{1\\6\\3} - \bMat{7\\1\\3} \\
            &= \bMat{-5\\11\\3}
        \end{align*}
        Finally, we get that
        \[
            A = \bMat{ 16 & -5 \\ -27 & 11 \\ -6 & 3}
        \]
    \end{solution}
    \emph{Hint:} Your final answer will have some large numbers, but you should not have to row reduce
    a matrix that looks too disgusting.
    \newpage
    \section*{Challenge Question}
    \question[20] Let $k,n$ be natural numbers\footnote{The natural numbers are integers greater than 0,
    in other words: $\setBasic{1,2,...}$} such that $k>n$. Let $\vec{v}_1,\dots,\vec{v}_k\in\R^n$.
    Argue why this list must be linearly dependent.
    \begin{solution}
        If we construct the augmented matrix
        \[
            \aMat{ccc|c}{\vec{v}_1 & \dots & \vec{v}_k & \vec{0}}
        \]
        then we have a $n\times k+1$ matrix. Since the final column is a column of $0s$ we will never get
        a pivot in this final column. We note that can have at most $n<k$ pivots
        as we can only have $1$ per row. Thus, we will have at least $k-n>0$ free
        variables. Therefore, the system
        \[
            a_1\vec{v}_1 + \cdots + a_k\vec{v}_k = \vec{0}
        \]
        will have infinitely many solutions, which means $\vec{v}_1,\dots,\vec{v}_k$ are linearly dependent as desired.
    \end{solution}
    \newpage
    \section*{Bonus Question}
    \question[10] Answer one of the following questions for 10 points. You may answer both, but I will
    only give you points for one!
    \begin{enumerate}
        \item Define the slang term ``glow up'' particularly in the context of:
            ``This person had a glow up.''
        \item Draw a picture of your favorite animal or character: either original or in media
    \end{enumerate}
\end{questions}

\end{document}
